% ===== §26 =====
\section{Emergence-Moments and the Outer Symbolic Order}
\label{sec:negotiation}

\noindent
The phenomenology closes at the edge where the couple's culture meets what lies outside it. Two things are treated here, and they are joined: first, the interpretive structure of the emergence-moment, the instant in which a private form is born and given its meaning; and second, the negotiation between the couple's endogenous culture and the outer symbolic order that surrounds and presses upon it. The first completes the account of how intimate culture is interpreted from within; the second opens the passage to the final part, for it is at this edge that the culture becomes vulnerable to the capture that the political economy of the last movement will analyze. The particular, having been developed in its irreducibility, is here returned to the general, and the return is not a retreat from the intimate case but the discovery, within it, of the point at which the outer world enters.

\subsection*{The co-signification of the emergence-moment}

The moment in which a private form emerges, described earlier as a moment of mutual self-encoding, has an interpretive structure that deserves to be drawn out, because it is where the couple's authority over their own meaning is most concentrated and most exposed. When a contingent occurrence is lifted into a form---when a mistaken word, a chance gesture, a shared perception is taken up and kept---the lifting is an act of \emph{co-signification}: the two, together, confer meaning on the occurrence, decide, usually without saying so, that this shall mean something and shall be kept. The co-signification is genuinely joint; it cannot be performed by one alone, for a form one party declared meaningful and the other did not take up would not be a shared form but a private wish. It requires the simultaneous assent of both, each reading in the other the recognition that the occurrence is to be kept, and this double assent is the couple exercising, at its purest, the authority the vacancy of the big Other left to them. In the emergence-moment the couple are their own big Other: they, and no external order, confer the meaning, and the fact that they can do this---that two people can, between them, make a thing mean---is the endogeneity of intimate culture caught at its source. This is the couple's freedom at its most concentrated, and, for the same reason, the point at which the outer order, if it can insert itself into the co-signification, gains the deepest purchase on the culture.

\subsection*{Negotiation with the outer symbolic order}

For the couple's culture does not emerge in a vacuum, and the endogeneity the paper has insisted on is never absolute. Around every intimate relation stands the outer symbolic order---the vast public stock of meanings, scripts, images, and expectations concerning what love is, how it is expressed, what a couple does, what a bond should look like---and the couple's culture is generated not in isolation from this order but in continual negotiation with it. The private forms a couple makes are made against a background of public forms they did not make: the culture supplies ready-made scripts for romance, for the marking of anniversaries, for the gestures of devotion, for the shape a shared life should take, and these press upon the couple's own generation, offering themselves as the forms the couple might adopt in place of generating their own. Every intimate culture is thus a mixture, in some proportion, of the endogenously generated and the externally supplied, and the proportion is not fixed but negotiated, continually, as the couple either takes up the public script or makes something of its own in its place. This negotiation is not in itself a corruption; a couple may take up a public form and make it genuinely theirs, charging the borrowed script with their own interwoven desire until it becomes, whatever its origin, a real part of their endogenous culture. The negotiation becomes corruption only when the public form is not appropriated but substituted---when the couple ceases to generate its own meaning and merely consumes the meaning the outer order supplies, so that the co-signification that should have been theirs is performed, in advance and elsewhere, by the market that made the script.

\subsection*{The doorway to capture}

It is here, precisely, that the phenomenology hands over to the political economy of the final part, and naming the handover completes the second movement. The point of vulnerability is the co-signification. Because the couple's meaning is conferred in the joint act of co-signification, and because that act can be colonized---its ready-made materials supplied, its outcomes pre-scripted, its very sense of what is worth keeping shaped in advance---the outer order has a route into the heart of the endogenous culture. It does not capture the culture by force, from outside; it captures it by supplying the forms in which the couple's own co-signification is carried out, so that the couple, generating what they take to be their own culture, are in fact re-enacting a script written elsewhere and owned by another. This is the mechanism the final part will anatomize as the capture of intimate culture by capital: not a seizure of a finished culture but a colonization of its generation, an insertion into the co-signifying moment of materials the market supplies and from which it draws its return. The second misrecognition of the previous section---the treatment of the bond's meaning as an exhibitable, convertible capital---is the subjective face of this capture, the disposition the outer order cultivates in the partners so that they will carry their own meaning out to the market; and the negotiation just described is the objective structure through which the capture operates.

\subsection*{The particular returned to the general}

With this the phenomenology of the dyad reaches its limit and returns the particular to the general from which it descended. The intimate culture, developed through this part in its irreducible specificity---its vacant big Other, its interwoven drive, its accelerated sedimentation, its intense registers, its fragility, its participatory hermeneutics---is shown, at its edge, to be embedded in and negotiated with an outer symbolic order that is itself, as the first part established, produced under material conditions and traversed by power and capital. The couple's private world is not a sanctuary outside the general history of culture but a particular site within it, connected to the whole at the point of negotiation, and therefore subject, at that point, to the forces the general account named. The third part takes up those forces directly. It asks how culture, including the intimate culture just described, is produced under the forces and relations of production; how capital, having exhausted its older frontiers, reaches through the negotiation just described into the couple's generation itself to capture the value it produces; and on what, in the end, the just generation of an intimate culture depends. The phenomenology has shown what intimate culture is and how it lives and dies; the political economy and ethics must now show how it is produced, how it is captured, and how it might be defended.
